% 迁移来源：NSFC_2026_CIESM-STSR/extraTex/2.2.特色与创新.tex
% 保真压缩：方法与判据细节见研究内容。
\subsubsection{项目特色}

利用新疆山地—盆地差异，以参考资料内配对、CIESM 非配对应用和独立台站三类验证，检验条件重建、模式迁移及是否更接近真实观测，统一考察大尺度状态、地形及多变量关系。

围绕 P02“算法或数据集”方向，形成可独立验收的\textbf{TerraSR-XJ 约束投影模块}与\textbf{新疆 0.1$^{\circ}$ 逐日四变量试验数据产品}，接口、规格及验收方式见研究目标。

\subsubsection{创新点}

\paragraph{检验参考场能否提供可学习的细节}
针对细网格被直接视为真值的问题，区分独立模拟、观测与插值细节，按研究方法中的 MAE 改善及逐变量 MCID 检验选源（气温初选 ERA5-Land，其余初选 CMFD）。检验仅用训练、验证资料，规则在申请阶段写定。全部候选源未通过的变量单独报告，不以平均收益掩盖其重建缺乏支持。

\paragraph{区分地形条件与联合建模的贡献}
在已有多变量时空降尺度基础上，固定预测因子、偏差校正、容量与预算，对比独立和共享表征，辨明联合收益及负迁移随地形的变化。新疆强逆温、地形抬升降水、绿洲冷岛与积雪反照率反馈提供丰富条件；所得认识用于 CIESM 重建，无稳定增益时报告适用边界。

\paragraph{分析大尺度一致性与局地细节的关系}
分别报告大尺度校正的信号变化、投影残差与经验性地形约束的独立样本效果。TerraSR-XJ 处理非嵌套、边界部分覆盖及秩亏，不可行时按既定规则降级。比较约束对不合理输出、极端及变量共同变化的影响。数学一致性不等于独立科学发现，聚合、非负与单位关系不代表完整水量或能量守恒。

\paragraph{事先确定评价标准}
申请阶段明确主次指标与规则：空间结构用无量纲 MCID，逐变量 MAE 用物理单位 MCID，非退化按 MAE 检查。归一尺度、权重及验证资料导出的阈值在查看测试集前固定存档，后续仅依原规则复算，区分支持、不支持与证据不足，禁止事后选阈值。
